% Introduction to Part IX. Unnumbered text: no \label except the \phantomsection target.
\phantomsection\label{intro:IX}
The localization of Part~VIII is of radius type. The polymer expansion of Part~X needs volume type: a fixed small factor for every cube of the dependence region (\hyperref[intro:VIII]{introduction to Part~VIII}). Part~IX provides it, as Ba\l aban and Dimock do, by random walks with weakening parameters. A product paving of the torus into boxes $\Box$ of side between $M/4$ and $M$ is fixed, with one parameter $s_\Box$ per box, and $s=(s_\Box)$ ranges over $\mathbb D_\kappa=\{|s_\Box|\le e^\kappa\}$ (Section~\ref{sec:weakenable}).

\paragraph{What weakening buys.} Let $\kappa\ge1$, let $f(s)$ be analytic in the interior of $\mathbb D_\kappa$ and continuous on it, with $|f|\le F$, and suppose that it decouples: if $s_\Box=0$ outside a union of boxes, then $f$ depends only on the parameters and the data near the component that contains a given polymer $Y$. Lemma~\ref{lem:localization} then writes $f(1)$ as a sum over polymers $Z\supset Y$ of pieces that depend on the data only near $Z$ and are bounded by $Fe^{-(\kappa-1)|Z\setminus Y|}$: a factor $e^{-(\kappa-1)}$ for every box. The proof has two steps. Write $f(1)$ by $g(1)=g(0)+\int_0^1g'(t)\,dt$ in each variable $s_\Box$ outside $Y$; the terms in which a differentiated box is not connected to $Y$ vanish by the decoupling. Each derivative is bounded by Cauchy's formula on a circle of radius $e^\kappa-1$, which stays in $\mathbb D_\kappa$; this gives the factor per box. So volume-type localization reduces to one task: to build every ingredient of the step as a decoupling function of $s$, analytic and bounded uniformly on $\mathbb D_\kappa$. Part~IX does this, and Part~X applies Lemma~\ref{lem:localization} (Section~\ref{sec:step-quantities}).

\paragraph{What this Part proves.}
\begin{itemize}
\item \emph{Walks in $L^2$} (Sections~\ref{sec:weakenable} and~\ref{sec:nested-walks}). A weakenable family is a sum over polymers $X_\omega$ of monomials $s^{n_\omega}$ times operators that depend on the data only near $X_\omega$ (Definition~\ref{def:weak-family}). On every domain, the inverse of an operator with local propagators and a factorized non-local part is such a family, for $M$ large, built by a walk in which every step carries the monomial of the boxes near its local domain (Theorem~\ref{thm:weak-inverse}). The walk runs at five nested levels (L1)--(L5), from the scalar operators to the fluctuation operator. Its outputs are $C_k^{1/2}(s)$ and the determinant densities (Corollary~\ref{cor:weak-covariance}), and the one-loop source as normalized graded activities with volume decay (Proposition~\ref{prop:one-loop-polymer}). The quadratic part of this source is marginal in four dimensions and must be extracted as a coupling, not bounded (Remark~\ref{rem:one-loop-polymer}(b)).
\item \emph{Pointwise bounds} (Sections~\ref{sec:pointwise-weak}--\ref{sec:tiers}). The background map needs pointwise $C^{1,\beta}$ bounds for $\bH(s)$, $\bC(s)$ and $\bC(s)D^T$, uniformly on $\mathbb D_\kappa$, with $\beta=\min(\alpha,\frac12)$. Interior estimates with $L^p$ data (Theorem~\ref{thm:lp-local}, Corollary~\ref{cor:lp-dirichlet}) and regularity tiers give them (Theorems~\ref{thm:tier-walk} and~\ref{thm:HC-pointwise}). In four dimensions $L^2$ data lie at the endpoint $p=\frac d2$, so the tiers run $L^2\to L^4\to L^8\to C^{1,\beta}$ (Lemma~\ref{lem:endpoint-factors}). No regularity at the Dirichlet cuts is used.
\item \emph{The weakened background map} (Section~\ref{sec:weak-background}). $A_*(s;B)$ is the fixed point of a weakened contraction. It is bounded uniformly on $\mathbb D_\kappa$, analytic in $(s,B)$, equal to $A_*(B)$ at $s=1$ at critical references, and it factorizes when $s_\Box=0$ outside a union of boxes (Theorem~\ref{thm:weak-bg}).
\end{itemize}

\paragraph{Imports and exports.} Part~IX builds on the Dirichlet operators, locality and complex perturbations of Part~III, on Propositions~\ref{prop:dirichlet-averaged} and~\ref{prop:local-fluct}, Theorem~\ref{thm:ID}(a)--(c), Section~\ref{sec:interior} and Lemma~\ref{lem:quadratic-bound} of Part~VIII, and on the contraction of Section~\ref{sec:var-solution}. Its clauses on complex references (in Propositions~\ref{prop:one-loop-polymer} and~\ref{prop:nl-tiers} and Theorem~\ref{thm:tier-walk}) use Proposition~\ref{prop:complex-ext}, and Section~\ref{subsec:torus-independence} uses Parts~X and~XI; both are best read with Part~X. Corollary~\ref{cor:weak-covariance}, Proposition~\ref{prop:one-loop-polymer} and Theorem~\ref{thm:weak-bg} are the inputs of the step in Part~X (Theorem~\ref{thm:sf-step}), and the torus independence of the one-loop pieces (Proposition~\ref{prop:one-loop-local}) is used in Parts~X and~XI. Theorem~\ref{thm:weak-inverse} also completes the gauge-field walk of Theorem~\ref{thm:random-walk}(b). Off the main line: Section~\ref{sec:pointwise-weak} is a discussion, whose Route~1 (boundary regularity) is not carried out, and Proposition~\ref{prop:schauder} is not used later.

\paragraph{Details for this Part.} The companion, chapter on Part~IX, begins with the main line through this Part and the $\delta$-budget of the walks, and then expands each of its steps, from the weakenable families to the weakened background map. Of Section~\ref{sec:pointwise-weak} it checks only the chain of exponents.
